\chapter{Satélites y movimiento planetario}\label{ch:g12:satellites-kepler}

La antena parabólica del balcón está atornillada y, sin embargo, se pasa el día
mirando a un punto fijo del cielo, \qty{36000}{km} más arriba; la estación espacial
cruza el cielo del atardecer en unos minutos; una constelación de relojes voladores le
dice a tu teléfono dónde estás. Ninguna de estas máquinas lleva un motor encendido:
todas se limitan a \emph{caer}. Este capítulo calcula a qué rapidez y a qué altura
debe caer cada una --- y después pesa planetas con la misma ley.

\section{Caer alrededor de la Tierra}

Todo descansa sobre la ley medida antes
(\cref{ch:g10:universal-gravitation}): la Tierra, de masa $M$, atrae a un \omterm{def:g10:universal-gravitation:satellite}{satélite} de
masa $m$ cuyo centro dista $r$ del suyo con una \omterm{def:g10:inertia:force}{fuerza} dirigida hacia el centro de la
Tierra y de módulo
\[ F = G\,\frac{m M}{r^2}, \qquad G = \qty{6.67e-11}{N.m^2/kg^2}. \]
Ojo con lo que es $r$: la distancia \emph{entre centros}. Un \omterm{def:g10:universal-gravitation:satellite}{satélite} a la altura $h$
tiene $r = R + h$, con $R = \qty{6.37e6}{m}$ el radio terrestre --- olvidarse de la
$R$ es el error clásico de este capítulo.

\begin{remark}[El cañón de Newton, revisitado]\label{rem:g12:satellites-kepler:cannon}
Disparada horizontalmente, una bala de cañón cae al suelo; más deprisa, cae más lejos;
a la rapidez justa, el suelo se curva y se aparta exactamente tan rápido como la bala
cae, y la caída se cierra sobre sí misma: una \omterm{def:g10:universal-gravitation:satellite}{órbita}. Un \omterm{def:g10:universal-gravitation:satellite}{satélite} no necesita motor
--- solo la rapidez \emph{lateral} adecuada.
\end{remark}

\begin{omfigure}
\begin{tikzpicture}[scale=1.0, font=\small]
  \fill[omExo!15] (0,0) circle (1.5); \draw[omExo, thick] (0,0) circle (1.5);
  \fill[omExo!60] (-0.22,1.44) -- (0,1.85) -- (0.22,1.44) -- cycle;
  \coordinate (L) at (0,1.85);
  \draw[omProp, thick] (L) .. controls (0.7,1.85) .. (50:1.5);
  \draw[omProp, thick] (L) .. controls (1.4,1.85) and (2.0,0.9) .. (0:1.5);
  \draw[-Stealth, omDef, thick] (0,1.85) arc (90:-260:1.85);
  \draw[-Stealth, omThm, thick] (L) .. controls (1.7,2.0) and (3.2,2.2)
    .. (4.3,2.9);
  \node[omDef, anchor=east] at (-1.85,1.1) {órbita};
  \node[omThm, anchor=west] at (3.35,2.35) {escapa};
  \node[omProp, anchor=west] at (1.35,0.55) {cae};
\end{tikzpicture}

{\small El cañón de Newton: disparada cada vez más deprisa, la bala cae más lejos
(naranja); hacia los \qty{7.9}{km/s} su caída iguala la curvatura de la Tierra y entra
en \omterm{def:g10:universal-gravitation:satellite}{órbita} (azul); aún más deprisa, escapa (rojo).}
\end{omfigure}

\section{La órbita circular: una rapidez, un periodo}

\begin{theorem}[Rapidez orbital]\label{thm:g12:satellites-kepler:speed}
Un \omterm{def:g10:universal-gravitation:satellite}{satélite} en \omterm{def:g10:universal-gravitation:satellite}{órbita} \omterm{def:g12:kinematics-2d:ucm}{circular uniforme} de radio $r$ alrededor de un cuerpo de masa
$M$ recorre su \omterm{def:g10:relative-motion:trajectory}{trayectoria} a la rapidez
\[ v = \sqrt{\frac{G M}{r}}, \]
que no depende de la masa del \omterm{def:g10:universal-gravitation:satellite}{satélite}.
\end{theorem}

\begin{proof}
La gravedad es la única \omterm{def:g10:inertia:force}{fuerza}, así que la \omterm{thm:g12:newtons-laws:second}{segunda ley de Newton}
(\cref{thm:g12:newtons-laws:second}) se lee $m\vect a = \vect F$; proyectémosla en la
\omterm{def:g12:kinematics-2d:frenet}{base de Frenet} (\cref{def:g12:kinematics-2d:frenet}). La \omterm{def:g10:inertia:force}{fuerza} apunta al centro: su
componente tangencial es nula, luego la rapidez es constante; su componente normal es
todo el $GmM/r^2$, y el movimiento circular tiene $a_n = v^2/r$
(\cref{thm:g12:kinematics-2d:centripetal}). Por tanto $m v^2/r = GmM/r^2$ y ---
dividiendo por $m$ --- $v^2 = GM/r$.
\end{proof}

\begin{omfigure}
\begin{tikzpicture}[scale=1.0, font=\small]
  \fill[omExo!15] (0,0) circle (0.85); \draw[omExo, thick] (0,0) circle (0.85);
  \node at (0,0) {Tierra}; \draw[dashed, omExo] (0,0) circle (2.3);
  \draw[thin] (0,0) -- (40:2.3) node[midway, above left] {$r$};
  \fill (40:2.3) circle (2.2pt) node[above right] {$S$};
  \draw[-Stealth, omThm, very thick] (40:2.3) -- (40:1.35)
    node[pos=0.85, below right] {$\vect F$};
  \draw[-Stealth, omDef, thick] (40:2.3) -- ++(130:1.15)
    node[above] {$\vect v$};
\end{tikzpicture}

{\small Todo el mecanismo de una \omterm{def:g10:universal-gravitation:satellite}{órbita}: velocidad tangente, \omterm{def:g10:inertia:force}{fuerza} hacia el centro
--- la gravedad nunca frena al \omterm{def:g10:universal-gravitation:satellite}{satélite}, solo curva su camino, para siempre.}
\end{omfigure}

\begin{proposition}[Periodo orbital]\label{prop:g12:satellites-kepler:period}
El \omterm{def:g10:universal-gravitation:satellite}{satélite} completa su \omterm{def:g10:universal-gravitation:satellite}{órbita} en el \omterm{def:g12:kinematics-2d:ucm}{periodo}
\[ T = \frac{2\pi r}{v} = 2\pi \sqrt{\frac{r^3}{G M}}. \]
\end{proposition}

\begin{proof}
Una vuelta es la circunferencia $2\pi r$ a rapidez constante; sustituyendo
$v = \sqrt{GM/r}$: $T = 2\pi r\sqrt{r/(GM)} = 2\pi\sqrt{r^3/(GM)}$.
\end{proof}

Ni $v$ ni $T$ contienen $m$: un tornillo que se suelta de la estación espacial orbita
justo a su lado. Y ambos decrecen con $r$: \emph{cuanto más alta es la \omterm{def:g10:universal-gravitation:satellite}{órbita}, más
lento va el \omterm{def:g10:universal-gravitation:satellite}{satélite}}, en rapidez ($v \propto 1/\sqrt r$) y en ángulo
($T \propto r^{3/2}$).

\begin{example}[La estación espacial]\label{ex:g12:satellites-kepler:iss}
La Estación Espacial Internacional vuela a la altura $h \approx \qty{400}{km}$, luego
$r = \num{6.37e6} + \num{4.0e5} = \qty{6.77e6}{m}$; con
$GM = \num{6.67e-11} \times \num{5.97e24} = \qty{3.98e14}{m^3/s^2}$,
$v = \sqrt{\num{3.98e14}/\num{6.77e6}} \approx \qty{7.7e3}{m/s}$ y
$T = 2\pi r/v \approx \qty{5.5e3}{s} \approx \qty{92}{min}$: una vuelta al planeta cada
hora y media, y unos dieciséis amaneceres al día para los astronautas.
\end{example}

\begin{example}[La Luna cae como la manzana]\label{ex:g12:satellites-kepler:moon}
La Luna, a $r = \qty{3.84e8}{m}$, orbita en $T = \qty{27.3}{d} = \qty{2.36e6}{s}$,
luego $v = 2\pi r/T \approx \qty{1.02e3}{m/s}$ y
$a = v^2/r \approx \qty{2.7e-3}{m/s^2}$ --- exactamente $g/60^2$, y la Luna está a
$60$ radios terrestres: el tirón que hace caer la manzana curva la Luna, debilitado
por el inverso del cuadrado. Cada segundo la Luna ``cae'' alrededor de \qty{1.4}{mm}
hacia nosotros, y su kilómetro lateral se la lleva alrededor --- la comprobación que
hizo Newton en 1666.
\end{example}

\section{Las leyes de Kepler}

Newton no adivinó el inverso del cuadrado: lo extrajo de tres regularidades que
Johannes Kepler destiló, entre 1609 y 1619, a partir de veinte años de posiciones
planetarias a simple vista tomadas por Tycho Brahe --- los mejores datos del mundo
anterior al telescopio.

\begin{definition}[Elipse]\label{def:g12:satellites-kepler:ellipse}
Una \emph{elipse}\index{elipse} es el conjunto de los puntos cuyas distancias a dos
puntos fijos, los \emph{focos}\index{foco (de una elipse)}, tienen suma constante; la
mitad de su diámetro mayor es el \emph{semieje mayor}\index{semieje mayor} $a$, y la
circunferencia es el caso en que los \omterm{def:g11:lenses-and-eye:focal}{focos} se confunden. Los puntos de la \omterm{def:g10:universal-gravitation:satellite}{órbita} de un
planeta más cercano y más lejano del Sol son su \emph{perihelio}\index{perihelio} y su
\emph{afelio}\index{afelio}.
\end{definition}

\begin{theorem}[Leyes de Kepler]\label{thm:g12:satellites-kepler:kepler}
\index{leyes de Kepler}%
Para los planetas que orbitan el Sol (de masa $M$):
\begin{enumerate}
  \item cada \omterm{def:g10:universal-gravitation:satellite}{órbita} es una \omterm{def:g12:satellites-kepler:ellipse}{elipse} con el Sol en uno de sus \omterm{def:g12:satellites-kepler:ellipse}{focos};
  \item el segmento Sol--planeta barre áreas iguales en tiempos iguales;
  \item el cociente entre el cuadrado del \omterm{def:g12:kinematics-2d:ucm}{periodo} y el cubo del \omterm{def:g12:satellites-kepler:ellipse}{semieje mayor} es el
        mismo para todos los planetas: $T^2/a^3 = 4\pi^2/(G M)$.
\end{enumerate}
Las mismas leyes gobiernan cualquier familia de \omterm{def:g10:universal-gravitation:satellite}{satélites} en torno a cualquier cuerpo
central, siendo $M$ la masa central. \admitted
\end{theorem}

\begin{remark}[Lo que podemos demostrar este curso]\label{rem:g12:satellites-kepler:proof}
Para una \omterm{def:g10:universal-gravitation:satellite}{órbita} \emph{circular} la tercera ley es nuestra: elevando al cuadrado la
\cref{prop:g12:satellites-kepler:period} sale $T^2/r^3 = 4\pi^2/(GM)$ --- una única
constante para todos los \omterm{def:g10:universal-gravitation:satellite}{satélites} del mismo centro. Que las \omterm{def:g12:satellites-kepler:ellipse}{elipses}, las áreas
iguales y el $a$ general se siguen del inverso del cuadrado se demuestra, con más
cálculo, en el volumen del primer año de universidad. Mientras tanto, la segunda ley
puede leerse: el planeta va más rápido en el \omterm{def:g12:satellites-kepler:ellipse}{perihelio}, donde el corto segmento que
barre tiene que darse prisa.
\end{remark}

\begin{omfigure}
\begin{tikzpicture}[scale=0.85, font=\small,
    declare function={re(\t)=2.025/(1+0.5*cos(\t));}]
  \draw[thick] plot[domain=0:360, samples=120, smooth]
    ({1.35+re(\x)*cos(\x)}, {re(\x)*sin(\x)});
  \fill[omDef!30] (1.35,0) -- plot[domain=-40:40, samples=25]
    ({1.35+re(\x)*cos(\x)}, {re(\x)*sin(\x)}) -- cycle;
  \fill[omDef!30] (1.35,0) -- plot[domain=175:185, samples=10]
    ({1.35+re(\x)*cos(\x)}, {re(\x)*sin(\x)}) -- cycle;
  \fill[omThm] (1.35,0) circle (2.6pt) node[below=2pt] {Sol};
  \draw[-Stealth, omDef, thick]
    ({1.35+re(40)*cos(40)}, {re(40)*sin(40)}) -- ++(-0.55,0.42);
  \draw[-Stealth, omDef, thick]
    ({1.35+re(175)*cos(175)}, {re(175)*sin(175)}) -- ++(-0.12,-0.5);
  \node[right] at (2.75,0) {$P$}; \node[left] at (-2.72,0) {$A$};
\end{tikzpicture}

{\small La segunda ley de Kepler: en un mes cerca del \omterm{def:g12:satellites-kepler:ellipse}{perihelio} $P$ el planeta barre un
sector corto y ancho; cerca del \omterm{def:g12:satellites-kepler:ellipse}{afelio} $A$, uno largo y estrecho --- áreas iguales: se
apresura cuando está cerca y remolonea cuando está lejos.}
\end{omfigure}

\begin{omfigure}
\begin{tikzpicture}
\begin{axis}[omaxis, width=8.8cm, height=5.8cm,
    xmode=log, ymode=log,
    xlabel={semieje mayor $a$ (unidades astronómicas)},
    ylabel={periodo $T$ (años)},
    xmin=0.25, xmax=50, ymin=0.15, ymax=400]
  \addplot[omThm, thick] coordinates {(0.3,0.1643) (40,253)};
  \addplot[only marks, mark=*, mark size=1.7pt, omDef] coordinates
    {(0.387,0.241) (0.723,0.615) (1,1) (1.524,1.881) (5.203,11.86)
     (9.537,29.46) (19.19,84.02) (30.07,164.8)};
  \node[below right, font=\scriptsize] at (axis cs:0.387,0.241) {Mercurio};
  \node[below right, font=\scriptsize] at (axis cs:1,1) {Tierra};
  \node[below right, font=\scriptsize] at (axis cs:5.203,11.86) {Júpiter};
  \node[below right, font=\scriptsize] at (axis cs:30.07,164.8) {Neptuno};
\end{axis}
\end{tikzpicture}

{\small Los ocho planetas en ejes logarítmicos: una recta de pendiente $3/2$, es decir,
$T^2 \propto a^3$ --- la tercera ley de Kepler, que en estas \omterm{def:g10:orders-of-magnitude:unit}{unidades} (\omterm{def:g10:orders-of-magnitude:lightyear}{unidades
astronómicas}, años) se lee sencillamente $T^2 = a^3$.}
\end{omfigure}

\section{Las órbitas en acción}

\begin{definition}[Órbita geoestacionaria]\label{def:g12:satellites-kepler:geo}
Un \omterm{def:g10:universal-gravitation:satellite}{satélite} es \emph{geoestacionario}\index{órbita geoestacionaria} cuando queda
colgado sobre un punto fijo del suelo. Su \omterm{def:g10:universal-gravitation:satellite}{órbita} debe ser circular, ecuatorial y de
\omterm{def:g12:kinematics-2d:ucm}{periodo} un \emph{día sidéreo}\index{día sidéreo}, $T = \qty{86164}{s}$ (\qty{23}{h}
\qty{56}{min} \qty{4}{s}): el tiempo que tarda la Tierra en dar una vuelta
\emph{respecto de las estrellas}. (El día de \qty{24}{h} es respecto del Sol, hacia el
cual la Tierra también avanza un grado al día.)
\end{definition}

\begin{example}[La altura geoestacionaria]\label{ex:g12:satellites-kepler:geoalt}
La tercera ley de Kepler, leída al revés, dicta el radio:
$r = \left(GMT^2/4\pi^2\right)^{1/3} =
\left(\num{3.98e14} \times (\num{86164})^2/4\pi^2\right)^{1/3} \approx
\qty{4.22e7}{m}$ --- una altura $h = r - R \approx \qty{3.58e7}{m} \approx
\qty{35800}{km}$, recorrida a $v = 2\pi r/T \approx \qty{3.1}{km/s}$. Todo repetidor
que ``se queda quieto'' se sitúa en esta única circunferencia sobre el ecuador --- no
hay otra.
\end{example}

\begin{example}[La constelación de navegación]\label{ex:g12:satellites-kepler:gps}
Los sistemas de navegación por \omterm{def:g10:universal-gravitation:satellite}{satélite} hacen volar una treintena de \omterm{def:g10:universal-gravitation:satellite}{satélites} a
$r \approx \qty{2.66e7}{m}$ (altura $\approx \qty{20200}{km}$), donde la
\cref{prop:g12:satellites-kepler:period} da $T \approx \qty{4.32e4}{s}$: medio \omterm{def:g12:satellites-kepler:geo}{día
sidéreo}, así que cada \omterm{def:g10:universal-gravitation:satellite}{satélite} vuelve a recorrer su traza sobre el suelo un día tras
otro. Tu receptor compara los tiempos de llegada de sus \omterm{def:g10:signals-and-waves:signal}{señales} de reloj con unas
\omterm{def:g10:universal-gravitation:satellite}{órbitas} conocidas al \omterm{def:g10:orders-of-magnitude:unit}{metro}.
\end{example}

\begin{method}[Pesar un mundo a partir de su satélite]\label{met:g12:satellites-kepler:weigh}
\index{pesar un planeta}%
Para medir la masa de un planeta o de una estrella: busca cualquier \omterm{def:g10:universal-gravitation:satellite}{satélite} suyo (una
luna, una sonda, un planeta), mide el radio $r$ y el \omterm{def:g12:kinematics-2d:ucm}{periodo} $T$ de su \omterm{def:g10:universal-gravitation:satellite}{órbita} e
invierte la tercera ley:
\[ M = \frac{4\pi^2 r^3}{G\,T^2}. \]
Esto pesa únicamente el cuerpo \emph{central} --- la masa del \omterm{def:g10:universal-gravitation:satellite}{satélite} se canceló en
el \cref{thm:g12:satellites-kepler:speed}. Todas las masas de la tabla de datos se han
obtenido así: la de la Tierra a partir de la Luna, la del Sol a partir de la propia
\omterm{def:g10:universal-gravitation:satellite}{órbita} terrestre.
\end{method}

\begin{remark}[Más alto es más lento --- pero cuesta más]\label{rem:g12:satellites-kepler:energy}
Como $v = \sqrt{GM/r}$, la Luna se arrastra a \qty{1}{km/s} mientras la estación corre
a \qty{7.7}{km/s}. Y sin embargo la \omterm{def:g10:universal-gravitation:satellite}{órbita} más alta es la más \emph{cara}: trepar
contra la gravedad almacena \omterm{def:g10:energy-conservation:energy}{energía} potencial (la cuenta se salda en el
\cref{ch:g12:work-and-energy}), y la ganancia en $E_p$ pesa más que la pérdida en
$E_k$. Por eso un cohete quema \emph{hacia adelante} para alcanzar una \omterm{def:g10:universal-gravitation:satellite}{órbita} más alta
y más lenta --- y un \omterm{def:g10:universal-gravitation:satellite}{satélite} frenado por el aire enrarecido cae en espiral \emph{hacia
abajo} y se acelera.
\end{remark}

\begin{remark}[Tabla de datos]\label{rem:g12:satellites-kepler:data}
Salvo indicación contraria: $G = \qty{6.67e-11}{N.m^2/kg^2}$; Tierra:
$M = \qty{5.97e24}{kg}$, $R = \qty{6.37e6}{m}$, \omterm{def:g12:satellites-kepler:geo}{día sidéreo} \qty{86164}{s}; \omterm{def:g10:universal-gravitation:satellite}{órbita} de
la Luna: $r = \qty{3.84e8}{m}$, $T = \qty{27.3}{d}$; Sol: $M = \qty{1.99e30}{kg}$;
distancia Tierra--Sol \qty{1.496e11}{m}; un año $= \qty{3.156e7}{s}$; Marte:
$M = \qty{6.42e23}{kg}$, $R = \qty{3.39e6}{m}$.
\end{remark}

\section{Ejercicios}

\begin{exercise}[$\star$]\label{exo:g12:satellites-kepler:1}
Un \omterm{def:g10:universal-gravitation:satellite}{satélite} de \qty{1000}{kg} está sobre la plataforma de lanzamiento y después vuela a
\qty{400}{km} de altura. Calcula la atracción terrestre en ambos lugares: ¿qué
porcentaje sobrevive allá arriba --- y por qué flotan sus ocupantes?
\end{exercise}

\begin{exercise}[$\star$]\label{exo:g12:satellites-kepler:2}
Para un \omterm{def:g10:universal-gravitation:satellite}{satélite} (teórico) que rozara la superficie terrestre ($r = R$), calcula la
rapidez orbital y el \omterm{def:g12:kinematics-2d:ucm}{periodo}. Compáralo con la bala de cañón de Newton
(\cref{rem:g12:satellites-kepler:cannon}).
\end{exercise}

\begin{exercise}[$\star$]\label{exo:g12:satellites-kepler:3}
Se traslada un \omterm{def:g10:universal-gravitation:satellite}{satélite} de una \omterm{def:g10:universal-gravitation:satellite}{órbita} circular de radio $r$ a otra de radio $2r$. ¿Por
qué factor cambian su rapidez y su \omterm{def:g12:kinematics-2d:ucm}{periodo}? ¿Dependen las respuestas de su masa?
\end{exercise}

\begin{exercise}[$\star$]\label{exo:g12:satellites-kepler:4}
Rehaz los números de la estación: a partir de $r = \qty{6.77e6}{m}$, calcula $v$, $T$ y
el número de \omterm{def:g10:universal-gravitation:satellite}{órbitas} completadas en \qty{24}{h}.
\end{exercise}

\begin{exercise}[$\star$]\label{exo:g12:satellites-kepler:5}
Enuncia las tres \omterm{thm:g12:satellites-kepler:kepler}{leyes de Kepler}. ¿En qué punto de su \omterm{def:g12:satellites-kepler:ellipse}{elipse} muy estirada va más
rápido un cometa, y qué ley lo dice? ¿Depende el \omterm{def:g12:kinematics-2d:ucm}{periodo} de un \omterm{def:g10:universal-gravitation:satellite}{satélite} de su propia
masa?
\end{exercise}

\begin{exercise}[$\star\star$]\label{exo:g12:satellites-kepler:6}
Pesa la Tierra: a partir del radio y del \omterm{def:g12:kinematics-2d:ucm}{periodo} de la \omterm{def:g10:universal-gravitation:satellite}{órbita} lunar (tabla de datos),
calcula la masa terrestre y compárala con el valor de la tabla.
\end{exercise}

\begin{exercise}[$\star\star$]\label{exo:g12:satellites-kepler:7}
Una empresa propone un \omterm{def:g10:universal-gravitation:satellite}{satélite} suspendido permanentemente sobre París (latitud
$49^\circ$ norte). A partir de la dirección de la gravedad, explica por qué toda
\omterm{def:g10:universal-gravitation:satellite}{órbita} circular está centrada en el centro de la Tierra; concluye que la propuesta es
imposible. ¿Dónde \emph{sí} puede quedarse colgado un \omterm{def:g10:universal-gravitation:satellite}{satélite}?
\end{exercise}

\begin{exercise}[$\star\star$]\label{exo:g12:satellites-kepler:8}
Fobos gira alrededor de Marte a $r = \qty{9.38e6}{m}$ en $T = \qty{7}{h}
\qty{39}{min}$. Deduce la masa de Marte; compárala con la tabla de datos.
\end{exercise}

\begin{exercise}[$\star\star$]\label{exo:g12:satellites-kepler:9}
Marte orbita el Sol con $a = \qty{2.279e11}{m}$ en \qty{687}{d}. Calcula $T^2/a^3$
para Marte y para la Tierra (tabla de datos), comprueba la tercera ley de Kepler y
deduce la masa del Sol.
\end{exercise}

\begin{exercise}[$\star\star$]\label{exo:g12:satellites-kepler:10}
La comprobación lunar de Newton: a partir de la tabla de datos, calcula la rapidez de
la Luna y su \omterm{thm:g12:kinematics-2d:centripetal}{aceleración centrípeta} $a = v^2/r$. Demuestra que es igual a $g/60^2$ y
di por qué eso convenció a Newton de que una sola ley gobierna la manzana y la Luna.
\end{exercise}

\begin{exercise}[$\star\star$]\label{exo:g12:satellites-kepler:11}
En \omterm{def:g10:orders-of-magnitude:unit}{unidades} solares (\omterm{def:g10:orders-of-magnitude:lightyear}{unidades astronómicas}, años) la tercera ley de Kepler se lee
$T^2 = a^3$: la recta de la gráfica logarítmica del curso. Sitúa en ella el asteroide
Ceres, $a = \qty{2.77}{au}$, y el cometa Halley, $a = \qty{17.8}{au}$: calcula ambos
\omterm{def:g12:kinematics-2d:ucm}{periodos}. Halley pasó por el \omterm{def:g12:satellites-kepler:ellipse}{perihelio} por última vez en 1986 --- ¿cuándo se le
espera de vuelta?
\end{exercise}

\begin{exercise}[$\star\star\star$]\label{exo:g12:satellites-kepler:12}
El aire enrarecido roza la estación y, sin embargo, su rapidez \emph{aumenta} a medida
que pierde altura. Calcula $v$ a las alturas de \qty{400}{km} y \qty{350}{km} y
resuelve después la paradoja: el \omterm{def:g10:inertia:inventory}{rozamiento} frena --- ¿de dónde sale la rapidez de
más?
\end{exercise}

\begin{exercise}[$\star\star\star$]\label{exo:g12:satellites-kepler:13}
Un telescopio descubre un exoplaneta que gira en torno a su estrella a
$r = \qty{7.48e9}{m}$ cada \qty{3.50}{d}. Pesa la estrella, en \omterm{def:g10:orders-of-magnitude:unit}{kilogramos} y en soles.
¿Qué masa --- la de la estrella o la del planeta --- \emph{no} puede obtenerse así, y
por qué?
\end{exercise}

\begin{exercise}[$\star\star\star$]\label{exo:g12:satellites-kepler:14}
Unos colonos marcianos quieren un repetidor ``areoestacionario'' sobre su base. Marte
gira respecto de las estrellas en \qty{24}{h} \qty{37}{min}: calcula el radio de la
\omterm{def:g10:universal-gravitation:satellite}{órbita} y la altura sobre la superficie marciana (tabla de datos).
\end{exercise}

\begin{exercise}[$\star\star\star$]\label{exo:g12:satellites-kepler:15}
El cometa Halley pasa por el \omterm{def:g12:satellites-kepler:ellipse}{perihelio} a $r_p = \qty{8.77e10}{m}$ moviéndose a
\qty{54.5}{km/s}; el \omterm{def:g12:satellites-kepler:ellipse}{afelio} está a $r_a = \qty{5.25e12}{m}$. En ambos extremos la
velocidad es perpendicular al segmento Sol--cometa, y la segunda ley de Kepler obliga
entonces a que $r_p v_p = r_a v_a$ (triángulos barridos iguales, de área
$\tfrac12 r v\,\Delta t$). Deduce la rapidez en el \omterm{def:g12:satellites-kepler:ellipse}{afelio} y comenta el cociente.
\end{exercise}

\section{Problema: Una plaza sobre el ecuador}

\begin{problem}[{Problema de fin de semana --- colocar un \omterm{def:g10:universal-gravitation:satellite}{satélite} sobre el ecuador: por
qué la única plaza de aparcamiento del cielo es una circunferencia a \qty{35800}{km} de
altura, cómo se compara con la estación que corre por debajo y cómo la misma ley pesa
Júpiter de camino}]
\label{pb:g12:satellites-kepler:1}
Una empresa de telecomunicaciones quiere un repetidor al que sus antenas de tierra,
atornilladas una sola vez, no tengan que perseguir nunca: un \omterm{def:g10:universal-gravitation:satellite}{satélite} \omterm{def:g12:satellites-kepler:geo}{geoestacionario}.
Tú eres la analista de la misión; usa la tabla de datos. Para la Parte~IV: Ío, la luna
de Júpiter, orbita a $r = \qty{4.22e8}{m}$ en $T = \qty{1.77}{d}$, y Júpiter gira
respecto de las estrellas en \qty{9}{h} \qty{56}{min}.

\medskip
\noindent\textbf{Parte I --- Las reglas para quedarse quieto.}
\begin{enumerate}
  \item El \omterm{def:g10:universal-gravitation:satellite}{satélite} debe permanecer sobre un punto fijo del suelo. ¿Cuál debe ser su
        \omterm{def:g12:kinematics-2d:ucm}{periodo} --- y por qué el \omterm{def:g12:satellites-kepler:geo}{día sidéreo} de \qty{86164}{s} y no \qty{24}{h}?
  \item En un movimiento circular la \omterm{def:g10:inertia:force}{fuerza} resultante apunta al centro de la
        circunferencia; aquí la única \omterm{def:g10:inertia:force}{fuerza} es la gravedad. Deduce que el centro de
        toda \omterm{def:g10:universal-gravitation:satellite}{órbita} circular es el centro de la Tierra.
  \item Una circunferencia suspendida sobre la latitud de París estaría centrada en el
        eje terrestre, al norte del centro. Concluye: ¿en qué plano debe estar una
        \omterm{def:g12:satellites-kepler:geo}{órbita geoestacionaria}?
  \item Escribe la \omterm{thm:g12:newtons-laws:second}{segunda ley de Newton} en la \omterm{def:g12:kinematics-2d:frenet}{base de Frenet} para una \omterm{def:g10:universal-gravitation:satellite}{órbita} circular
        de radio $r$ y deduce $v = \sqrt{GM/r}$.
  \item Deduce $T = 2\pi\sqrt{r^3/(GM)}$ y reescríbelo como la tercera ley de Kepler,
        $T^2/r^3 = 4\pi^2/(GM)$.
\end{enumerate}

\medskip
\noindent\textbf{Parte II --- La única circunferencia que sirve.}
\begin{enumerate}[resume]
  \item Despeja $r$ en la tercera ley y calcula el radio \omterm{def:g12:satellites-kepler:geo}{geoestacionario}.
  \item Deduce la altura sobre el suelo.
  \item Calcula la rapidez orbital a partir de $v = 2\pi r/T$.
  \item Compruébala con $v = \sqrt{GM/r}$.
  \item El \omterm{def:g10:universal-gravitation:satellite}{satélite} de la empresa tiene una masa de \qty{4.0e3}{kg}; el de una rival
        pesa el doble. Compara sus radios \omterm{def:g12:satellites-kepler:geo}{geoestacionarios} y di qué cambia la masa de
        más en el lanzamiento.
\end{enumerate}

\medskip
\noindent\textbf{Parte III --- La estación, más abajo.}
\begin{enumerate}[resume]
  \item La estación espacial vuela a $r = \qty{6.77e6}{m}$: calcula su rapidez.
  \item Calcula su \omterm{def:g12:kinematics-2d:ucm}{periodo}, en segundos y en minutos.
  \item ¿Cuántas \omterm{def:g10:universal-gravitation:satellite}{órbitas} completa en \qty{24}{h} --- aproximadamente cuántos amaneceres
        contemplan los astronautas al día?
  \item Calcula el cociente de radios $r_{\text{geo}}/r_{\text{EEI}}$ y comprueba que el
        cociente de \omterm{def:g12:kinematics-2d:ucm}{periodos} es ese número elevado a $3/2$, como exige la tercera ley
        de Kepler.
  \item ¿Qué \omterm{def:g10:universal-gravitation:satellite}{satélite} se mueve más deprisa, y cuál costó más \omterm{def:g10:energy-conservation:energy}{energía} por \omterm{def:g10:orders-of-magnitude:unit}{kilogramo}
        colocar? Reconcilia ambas cosas en una frase.
\end{enumerate}

\medskip
\noindent\textbf{Parte IV --- Pesar Júpiter de paso.}
\begin{enumerate}[resume]
  \item A partir de la tercera ley de Kepler aplicada a Ío, expresa la masa de Júpiter
        $M_J$ en función de $G$, $r$ y $T$.
  \item Calcula $M_J$.
  \item ¿A cuántas masas terrestres equivale? ¿Aproximadamente qué fracción de la masa
        del Sol?
  \item Un repetidor ``joviestacionario'' debería quedar colgado sobre un punto de las
        nubes de Júpiter: calcula el radio de su \omterm{def:g10:universal-gravitation:satellite}{órbita}; compáralo con el de Ío.
  \item Informe al consejo, en dos frases: la altura donde la empresa debe aparcar su
        repetidor y la masa de Júpiter que esa misma ley ha entregado gratis.
\end{enumerate}
\end{problem}
